\section{Difféomorphismes locaux et fonctions  implicites}
\subsection{Difféomorphisme locaux}
\newdef{Soient $E, F \subset \rn$ ouverts non-vides. Une application $\bm f : E \to F$ est un \emph{homéomorphisme} si elle est bijective et si, $\bm f$ et son inverse $\bm g$ sont tout deux continues. 
\\\\
De plus pour $\bm f : E \to F$ une application bijective de classe $C^k$, on dit que $\bm f$ un \emph{k-difféomorphisme} si son inverse $\bm g$ est de classe $C^k$. 
\\\\
La définition s'étend localement, on parlera alors de homéomorphisme local  et de $k$-difféomorphisme en $\bm x_0 \in E$
}{}

\newlem{Soit $\bm f : E \to F$ un difféomorphisme locale en tout point $\bm x_0 \in E$ alors si $\bm f : E \to F$ est une bijection entre $E$ et $F$, alors elle est un difféomorphisme globale.
}{} 

\subsection{Théorème d'inversion locale}

\newlem{Soit $\bm f : E \to F$ un difféomorphisme locale en $\bm x_0 \in E$. Alors $\det D\bm f(\bm x_0) \neq 0$.
}{}

\prf{Puisque $\bm f : E \to F$ est un difféomorphisme locale en $\bm x_0$, il existe un ouvert $U \subset E$ contenant $\bm x_0$ et un ouvert $V \subset F$ contenant $\bm f(\bm x_0) = \bm y_0$ tels que $\bm f : U \to V$ est bijective et la fonction $\bm g : V \to U$ est de classe $C^1$. Puisque $\bm f$ et $\bm g$ sont différentiable respectivement en $\bm x_0$ et $\bm y_0$ alors
	\[
	D \bm f(\bm x_0) D \bm g(\bm y_0) = D \bm g(\bm y_0) D \bm f(\bm x_0) = \mathrm I_n
	\]
	ce qui implique que $D \bm f(\bm x_0)$ est inversible et donc $\det D \bm f(\bm x_0) \neq 0$.
}

\newthm{Soient $K \subset \rn$ un fermé non-vide et $\bm \phi : K \to \rn$ une application telle que,
	\begin{enumerate}
		\item $\bm \phi(K) \subset K$
		\item $\exists \rho \in ]0,1[, ~\forall \bm w, \bm v \in K, \quad \norme{\bm v - \bm w} \leqslant \rho \norme{\bm v - \bm w}$
	\end{enumerate}
Alors il existe un \emph{unique} $\bm{\tilde v} \in K$ tel que $\bm \phi(\bm{\tilde v}) = \bm{\tilde v}$.
}{Théorème du point fixe de Banach}

\prf{On fixe $\bm v^0 \in K$ et on définit la suite définit par récurrence,
	\[
	\forall k \in \N, \quad \bm v^{k+1} = \bm \phi (\bm v^k)
	\]
	On remarque alors que pour tout $k \in \N$ on a,
	\[
	\norme{\bm v^{k+1} - \bm v^k} = \norme{\bm \phi(\bm v^k) - \bm \phi (\bm v^{k-1})} \leqslant \rho \norme{\bm v^k - \bm v^{k-1}} \leqslant \ldots \leqslant \rho^k \norme{\bm v^1 - \bm v^0}
	\]
	Ainsi pour $\varepsilon > 0$ et $k > m > 0$,
	\[
	\norme{\bm v^k - \bm v^m} \leqslant \sum_{i=m}^k \norme{\bm v^{i+1} - \bm v^i} \leqslant \sum_{i=m}^{k-1} \rho^i \norme{\bm v^1 - \bm v^0} \leqslant \frac{\rho^m}{1 - \rho} \norme{\bm v^1 - \bm v^0} \leqslant \varepsilon
	\]
	si $m$ est suffisamment grand. Ainsi la suite est de Cauchy, donc convergente. On pose donc $\bm{\tilde v}$ la limite de $(\bm v^k)$. Alors $\bm{\tilde v} \in K$ puisque $K$ est fermé, et 
	\begin{align*}
		\norme{\bm{\tilde v} - \bm \phi(\bm{\tilde v})} \leqslant \norme{\bm{\tilde v} - \bm v^{k+1}} &+ \norme{\bm v^{k+1} - \bm \phi(\bm{\tilde v})} = \norme{\bm{\tilde v} - \bm v^{k+1}} + \norme{\bm \phi (v^k) - \bm \phi(\bm{\tilde v})} \\
		& \leqslant \norme{\bm{\tilde v} - \bm v^{k+1}} + \rho \norme{\bm v^k - \bm{\tilde v}} \underset{k \to \infty}{\longrightarrow} 0
	\end{align*}
	D'où $\bm \phi(\bm{\tilde v}) = \bm{\tilde v}$. Supposons désormais qu'il existe un autre $\bm{\tilde u} \in K$ tel que $\bm \phi(\bm{\tilde u}) = \bm{\tilde u}$. Alors
	\[
	\norme{\bm{\tilde v} - \bm{\tilde u}} = \norme{\bm \phi(\bm{\tilde v}) - \bm \phi(\bm{\tilde u})} \leqslant \rho \norme{\bm{\tilde u} - \bm{\tilde v}}
	\]
	D'où $\bm{\tilde v} = \bm{\tilde u}$.
}

\newlem{Soit $A \in M_n(\R)$, on note $\|| A \||$ la \emph{norme spectrale} de $A$ définie par $\|| A \|| = \sup_{\norme{\bm \xi }=1} \norme{A \bm \xi}$, et par $\| A \|_F$ la \emph{norme de Frobenius} de $A$ définie par $\|A \|_F = \sqrt{\sum_{i,j = 1}^n A_{ij}^2}$. Alors $\|| A \|| \leqslant \|A \|_F$.
}{Norme spectrale}

\prf{Grâce à l'inégalité de Cauchy-Schwarz,
	\[
	\|| A \||^2 = \sup_{\bm \xi \in \rn,~\norme{\bm \xi }=1} \sum_{i=1}^n \left(\sum_{j=1}^n A_{ij} \xi_j\right)^2 \leqslant \sup_{\bm \xi \in \rn,~\norme{\bm \xi }=1} \sum_{i=1}^n  \left(\sum_{j=1}^n A_{ij}^2\right)\left(\sum_{j=1}^n \xi_j^2\right) = \norme{A}_F^2
	\]
}

\newlem{Soit $\bm f : [a,b] \to \rn$ une fonction continue, alors
	\[
	\left \| \int_a^b \bm f(t) ~\mathrm dt \right\| \leqslant \int_a^b \norme{\bm f(t)} ~\mathrm dt
	\]
}{Norme intégrale}

\prf{On note $\bm v \in \rn$ le vecteur de composantes $v_i = \int_a^b f_i(t) ~\mathrm dt$. Alors,
	\[
	\left\| \int_a^b \bm f(t) ~\mathrm dt \right\|^2 = \norme{\bm v}^2 = \sum_{i=1}^n v_i v_i = \int_a^b \sum_{i=1}^n v_i f_i(t) ~\mathrm dt \leqslant \int_a^b \norme{\bm v} \norme{\bm f(t)} ~\mathrm dt
	\]
}

\newthm{Soit $f : E \to F$ une fonction continue de classe $C^1$ et $\bm x_0 \in E$. Si $\det D \bm f (\bm x_0) \neq 0$, alors $\bm f$ est un difféomorphisme locale en $\bm x_0$. De plus $D \bm g(\bm f(\bm x)) = D \bm f(\bm x)^{-1}$ pour tout $\bm x \in U$.
}{Théorème d'inversion locale}

\prf{La preuve se décompose en trois étapes : 
	\begin{enumerate}
		\item On montre qu'il existe $r, \tilde r > 0$ tel que $\forall \bm y \in B(\bm y_0, \tilde r)$ l'équation $\bm f(\bm x) = \bm y$ pour $\bm x \in B(\bm x_0,r)$ possède une unique solution en utilisant le théorème du point fixe de Banach (\ref{Théorème du point fixe de Banach})
		
		\item On pose donc $V = B(\bm y_0, \tilde r)$ et $U = B(\bm x_0, r) \cap \bm f^{-1}(V)$, on montre $\bm f : U \to V$ est une bijection et la fonction inverse $\bm g : V \to U$ est continue.
		
		\item On montre finalement $\bm g: V \to U$ est de classe $C^1$ et $\forall \bm x \in U, \quad D \bm g(\bm f(\bm x)) = D \bm f(\bm x)^{-1}$.
	\end{enumerate}
	\textit{Première étape} : Si on considère l'application,
	\[
	\begin{array}{cccc}
		\varphi : & E & \to &F \\
		& \bm x & \mapsto & \mathrm I_n - D\bm f(\bm x_0)^{-1} D\bm f(\bm x)
	\end{array}
	\]
	on remarque alors que $\phi$ est continue et s'annule en $\bm x_0$. Ainsi il existe $r > 0$ tel que,
	\[
	\forall \bm x \in \overline B(\bm x_0, r), ~\forall i,j \in \entiers, \quad \left|(\mathrm I_n - D \bm f(\bm x_0)^{-1} D \bm f(\bm x))_{ij} \right| \leqslant \frac{1}{2n}
	\]
	et de plus donc par continuité de l'application $\bm x \mapsto \det D \bm f(\bm x)$, on en déduit que pour tout $\bm x \in \overline B(\bm x_0, r), ~\det D \bm f(\bm x) \neq 0$. De plus on remarque que l'équation $\bm f(\bm x) = \bm y$ est équivalente à $\bm x = \varphi (x)$ où $\bm \varphi (\bm x) = \bm x - D \bm f(\bm x_0)^{-1} (\bm f(\bm x) - \bm y)$. En particulier $\bm \phi \in C^1$. Alors $\bm f(\bm x) = \bm y$ possède une solution dans $\overline B(\bm x_0, r)$ si et seulement si $\bm \phi$ possède un point fixe dans $\overline B(\bm x_0, r)$. Mais pour $\bm x\in \overline B(\bm x_0, r)$, 
	\[
	D \bm \phi = \mathrm I_n - D \bm f(\bm x_0)^{-1} D \bm f(\bm x)
	\]
	et donc $\left| \frac{\partial \bm \phi_i}{\partial x_j}(\bm x) \right| \leqslant \frac{1}{2n}$. Ainsi on en déduit que,
	\[
	\forall \bm x \in \overline B(\bm x_0, r), \quad \|| D \bm \phi(\bm x) \|| \leqslant \norme{D \bm \phi(\bm x)}_F = \left( \sum_{i,j = 1}^n \left|\frac{\partial \bm \phi_i}{\partial x_j}(\bm x)\right|^2\right)^{\frac{1}{2}} \leqslant \frac{1}{2}
	\]
	Par le théorème des accroissement finis version intégrales \ref{Théorème des accroissements finis} et les lemmes (\ref{Norme spectrale}-\ref{Norme intégrale}) on a pour $\bm x_1, \bm x_2 \in \overline B(\bm x_0, r)$,
	
	\begin{align*}
		\norme{\bm \phi(\bm x_2) - \bm \phi(\bm x_1)} &= \| \int_0^1 D \bm \phi (\bm x_1 + t(\bm x_2 - \bm x_1)) (\bm x_2 - \bm x_1) ~\mathrm dt \\
													&\leqslant \int_0^1 \norme{D \bm \phi(\bm x_1 + t(\bm x_2 - \bm x_1)) (\bm x_2 - \bm x_1)} ~\mathrm dt \\
													& \leqslant \int_0^1 \|| D \bm \phi(\bm x_1 + t(\bm x_2 - \bm x_1)) \|| \norme{\bm x_2 - \bm x_1} ~\mathrm dt \leqslant \frac{1}{2} \norme{\bm x_2 - \bm x_1}
	\end{align*}
	Ainsi $\bm \phi$ est contractante sur $\overline B(\bm x_0, r)$ pour tout $\bm x \in \rn$. De plus, pour tout $\bm x \in \overline B(\bm x_0,r)$ et $\bm y \in B(\bm y_0, \tilde r)$, avec $\tilde r = \frac{r}{2 \|| D \bm f(\bm x_0)^{-1}\||}$, on a
	\begin{align*}
		\norme{\bm \phi(\bm x) - \bm  x_0}   &\leqslant \norme{\bm \phi(\bm x) - \bm \phi(\bm x_0)} + \norme{\bm \phi (\bm x_0) - \bm x_0} \\
												&\leqslant \frac{1}{2} \norme{\bm x -\bm x_0} + \norme{D \bm f(\bm x_0)^{-1} (\bm y_0 - \bm y)} \leqslant \frac{r}{2} + \|| D \bm f(\bm x_0)^{-1} \||  \norme{\bm y_0 - \bm y} < r		
	\end{align*}
	Donc si on prend $\bm y \in B(\bm y_0, \tilde r)$, on a $\bm \phi (\overline B(\bm x_0, r)) \subset B(\bm x_0, r)$. Par conséquent, $\bm \phi : \overline B(\bm x_0, r) \to \overline B(\bm x_0, r)$ possède un unique point fixe $\bm x \in \overline B(\bm x_0, r)$. Mais comme $\bm \phi (\overline B(\bm x_0, r)) \subset B(\bm x_0, r)$, ce point fixe est dans la boule ouverte $B(\bm x_0, r)$. On a donc montré que $\forall \bm y \in B(\bm y_0, \tilde r)$, il existe un unique $\bm x \in B(\bm x_0, r)$ tel que $\bm x = \bm \phi(\bm x)$, c'est-à-dire $\bm f(\bm x) = \bm y$. \\ \\
	\textit{Deuxième étape} : Soit $V = B(\bm y_0, \tilde r)$ et $U = B(\bm x_0, r) \cap \bm f^{-1} (V)$. On remarque que $\bm f^{-1} (V)$ est ouvert car la pré-image d'un ouvert par une fonction continue est ouvert, donc $U$ ouvert. On sait que, d'après la première étape que $\bm f: U \to V$ est une bijection, on peut donc définir la fonction inverse $\bm g : V \to U$. Soit $\varepsilon > 0$ et $\delta = \frac{\varepsilon}{2 \|| D \bm f(\bm x_0)^{-1} \||}$. Ainsi pour tout $\bm y_1, \bm y_2 \in V$ tels que $\norme{\bm y_1 - \bm y_2} \leqslant \delta$ et en notant $\bm x_1 = \bm g(\bm y_1)$ et $\bm x_2 = \bm g(\bm y_2)$, on a
	\begin{align*}
		\norme{\bm x_1 - \bm x_2} &= \norme{\bm \phi_{\bm y_1}(\bm x_1) - \bm \phi_{\bm y_2} (\bm x_2)} \\
									& \leqslant \norme{\bm \phi_{\bm y_1} (\bm x_1) - \bm \phi_{\bm y_1} (\bm x_2)} + \norme{\bm \phi_{\bm y_1} (\bm x_2) - \phi_{\bm y_2} (\bm x_2)} \\
									&\leqslant \frac{1}{2} \norme{\bm x_1 - \bm x_2} + \norme{D \bm f(\bm x_0)^{-1}(\bm y_1 - \bm y_2)}
	\end{align*}
	et donc,
	\[
	\norme{\bm x_1 - \bm x_2} = \norme{\bm g (\bm y_1) - \bm g(\bm y_2)} \leqslant 2 \|| D \bm f(\bm x_0)^{-1} \|| \norme{\bm y_1 - \bm y_2} \leqslant \varepsilon
	\]
	ce qui montre que $\bm g : V \to U$ est de Lipschitz, ce qui implique la continuité.
	\\\\
	\textit{Troisième étape} : Il reste à démontrer que $\bm g : V \to U$ est de classe $C^1$ et $D \bm g(\bm y) = D \bm f(\bm x)^{-1}$. Par le choix de $r$, $D \bm f(\bm x)$ est inversible pour tout $\bm x \in B(\bm x_0, r)$. Soit maintenant $\bm y_1 \in V$ et $\bm x = \bm g(\bm y_1)$. Puisque $\bm f$ est de classe $C^1$, on a
	\[
	\bm f(\bm x_1) - \bm f(\bm x) = D \bm f(\bm x) (\bm x_1 - \bm x) + \bm R_{\bm f}(\bm x_1)
	\]
	où $\bm R_{\bm f} (\bm x_1) = o(\norme{\bm x_1 - \bm x})$. Ceci implique
	\[
	\bm x_1 - \bm x = \bm g(\bm y_1) - \bm g(\bm y) = D\bm f(\bm x)^{-1}(\bm y_1 - \bm y) - \underbrace{D \bm f(\bm x)^{-1} \bm R_{\bm f}(\bm x_1)}_{\bm R_{\bm g}(\bm y_1)}
	\]
	Il nous reste à montrer que $\bm R_{\bm g}(\bm y_1) = o(\norme{\bm y_1 - \bm y})$,
	\[
	\lim_{\bm y_1 \to \bm y} \frac{\norme{\bm R_{\bm g}(\bm y_1)}}{\bm y_1 - \bm y} \leqslant \lim_{\bm y_1 \to \bm y} \|| D \bm f(\bm x)^{-1} \|| \frac{\norme{\bm R_{\bm f}(\bm x_1)}}{\norme{\bm y_1 - \bm y}} = \|| D \bm f(\bm x)^{-1} \|| \lim_{\bm y_1 \to \bm y} \frac{\norme{\bm R_{\bm f}(\bm x_1)}}{\norme{\bm x_1 - \bm x}} \frac{\norme{\bm x_1 - \bm x}}{\norme{\bm y_1 - \bm y}}
	\]
	Mais on a vu, lors de la seconde étape que,
	\[
	\norme{\bm x_1 - \bm x} \leqslant 2 \|| D \bm f(\bm x_0)^{-1} \|| \norme{\bm y_1 - \bm y}
	\]
	et donc,
	\[
	\lim_{\bm y_1 \to \bm y} \frac{\bm R_{\bm g}(\bm y_1)}{\norme{\bm y_1 - \bm y}} \leqslant 2 \|| D \bm f(\bm x)^{-1} \|| \|| D \bm f(\bm x_0)^{-1} \||  \lim_{\bm x_1 \to \bm x} \frac{\bm R_{\bm f}(\bm x_1)}{\norme{\bm x_1 - \bm x}} = 0
	\]
	On conclut que $\bm g : V \to U$ est différentiable en tout $\bm y \in V$ et $D \bm g(\bm y) = D \bm f(\bm x)^{-1}$, avec $\bm x = \bm g(\bm y)$. Enfin puisque $D \bm f(\bm x)$ en continue en tout $\bm x \in B(\bm x_0, r), ~ D \bm f(\bm x)^{-1}$ est continue pour vu que $\det D \bm f(\bm x) \neq 0$, ce qui est vrai dans $B(\bm x_0, r)$. Ceci montre que $\bm g$ est de classe $C^1$.
}


\subsection{Hypersurfaces et fonctions implicites}

\newdef{Soit $\Sigma \subset \R^{n+1}, ~\bm z_0 \in \Sigma$ et $k \in \N$. On dit que $\Sigma$ est une \emph{hypersurface} de classe $C^k$ autour $\bm z_0$ si elle est le graphe d'une fonction de classe $C^k$ localement autour de $\bm z_0$, c'est-à-dire s'il existe un voisinage $V$ de $\bm z_0$, un indice $i \in \llbracket 1, n+ 1 \rrbracket$, un ouvert $U \subset \rn$ et une fonction $\phi : U \to \R$ de classe $C^k$ tels que
	\[
	\Sigma \cap V = G(\phi) = \{ \bm x \in \R^{n+1} ~|~ x_i = \phi(\bm x_{i}), ~\bm x_{i} = (x_1, \ldots, x_{i-1}, x_{i+1}, \ldots, x_{n+1}) \in U \}
	\]
On considère donc des ensembles de la forme,
\[
\Sigma = \{ \bm E ~|~ f(\bm x) = 0\}
\]
avec $f : E \to \R$. Ainsi par ce qui précède, alors $\Sigma$ peut être représenté comme le graphe d'une fonction $\phi : U \to  \R$ localement autour de $\bm z_0$, alors
\[
\Sigma \cap V = G(\phi) = \{ \bm x \in \R^{n+1} ~|~ x_i = \phi(\bm x_i), ~ \bm x_i \in U\} = \{\bm x \in V ~|~ f(\bm x) = 0\}
\]
On dit alors dans ce cas que l'équation $f(\bm x) = 0$ définit \emph{implicitement} une fonction $x_i = \phi(\bm x_i)$ localement autour du point $\bm z_0$.
}{}

\subsection{Théorèmes des fonctions implicites}
\newthm{Soit $E \subset \R^{n+1}$ un ouvert non vide, $f : \rn \to \R$ une fonction de classe $C^1$, $\Sigma = \{(\bm x, y) \in E ~|~ f(\bm x, y) = 0\}$ et $\bm z_0 = (\bm x_0, y_0) \in \Sigma$ tel que $\frac{\partial f}{\partial y}(\bm x_0, y_0) \neq 0$. Alors il existe un ouvert $U = B(\bm x_0, \delta) \subset \rn$ de $\bm x_0$, un ouvert $V \subset E$ contenant $\bm z_0$ et une unique application $\phi : U \to \R$ de classe $C^1$ telle que : 
	\begin{enumerate}
		\item $y_0 = \phi(\bm x_0)$
		\item $\forall \bm x \in U, \quad (\bm x, \phi(\bm x)) \et f(\bm x, \phi(\bm x)) = 0$
		\item $G(\phi) = \Sigma \cap V$
	\end{enumerate}
	De plus, on a pour tout $\bm x \in U, \quad \frac{\partial f}{\partial y} (\bm x, \phi(\bm x)) \neq 0$ et,
	\[
	\frac{\partial \phi}{\partial x_i} (\bm x) = - \frac{\frac{\partial f}{\partial x_i} (\bm x, \phi(\bm x))}{\frac{\partial f}{\partial y}(\bm x, \phi(\bm x))}
	\]
	et si $f \in C^k$ alors $\phi \in C^k$.
}{Théorème des fonctions implicites}


\prf{Sans perte de généralité, supposons que $\frac{\partial f}{\partial y} (\bm x_0, y_0) > 0$. Puisque $\frac{\partial f}{\partial y}$ est continue sur $E$, il existe donc $\delta_1, \delta_2 > 0$ tels que, pour tout $(\bm x, y) \in B(\bm x_0, \delta_1) \times [y_0 - \delta_2, y_0 + \delta_2]$, on a
\[
\frac{\partial f}{\partial y} (\bm x, y) > 0
\]
Ainsi pour tout $\bm x \in B(\bm x_0, \delta_1)$, la fonction $y \mapsto g_{\bm x}(y) = f(\bm x, y)$ est strictement croissante dans l'intervalle $[y_0 - \delta_2, y_0 + \delta_2]$. En particulier,
\[
g_{\bm x_0} (y_0 - \delta_2)< g_{\bm x_0}(y_0) = 0 < g_{\bm x_0}(y_0 + \delta_2)m
\]
Mais puisque $f$ est continue, il existe $0 < \delta < \delta_1$ tel que,
\[
\forall \bm x \in U := B(\bm x_0, \delta), \quad g_{\bm x} (y - \delta_2) =  f(\bm x, y_0 - \delta_2) < 0 \et g_{\bm x}(y + \delta 2) = f(\bm x, y_0 + \delta_2) > 0
\]
Ainsi pour $\bm x \in U$, $g_{\bm x} (y)$ est continue est strictement croissante sur l'intervalle $[y_0 - \delta_2, y_0 + \delta_2]$, et donc il existe un unique $y = \phi(x) \in ]y_0 - \delta_2, y_0 + \delta_2[$ tel que $g_{\bm x}(y) = f(\bm x, y) = 0$. Cette procédure permet de définir de façon unique une fonction $\phi : U \to \R$ telle que,
\[
\phi(\bm x_0) = y_0 \et f(\bm x, \phi(\bm x)) = 0
\]
De plus si on note $V = B(\bm x_0, \delta) \times ]y_0 - \delta_2, y_0 + \delta_2[$, on a $G(\phi) = \Sigma \cap V$.\\\\
\textit{Continuité de $\phi$ : } Soit $\bm{\overline x} \in U$ et $\overline y = \phi(\bm{\overline x})$. Pour tout $\varepsilon \in ]0, \delta_2 - |\overline y - y_0|]$, on considère
\[
W_\varepsilon = B(\bm x_0, \delta_1) \times [\overline y - \varepsilon, \overline y + \varepsilon]
\]
Ainsi par construction $W_\varepsilon \subset W$, donc $\frac{\partial f}{\partial y} > 0$ sur $W_\varepsilon$. De la même manière que auparavant, mais sur $W_\varepsilon$ au lieu de $W$, il existe $\delta_\varepsilon >0$ et une fonction $\overline \phi : B(\bm{\overline x}, \delta_\varepsilon) \cap U : \to \R$ tels que 
\[
\overline y = \overline \phi(\bm{\overline x}) \et f(\bm x, \overline \phi(\bm x)) = 0
\]
Mais par unicité de la fonction implicite sur $W$ et le choix de $W_\varepsilon$, on a $\overline \phi(\bm x) = \phi(\bm x) \in [\overline y, \varepsilon, \overline y + \varepsilon]$ pour tout $\bm x \in B(\bm{\overline x}, \delta_\varepsilon)\cap U$. On a donc montré que pour tout $\varepsilon > 0$ il existe $\delta_\varepsilon > 0$ tel que,
\[
\forall \bm x \in B(\bm{\overline x} \cap U, \delta_\varepsilon), \quad |\phi(\bm x) - \phi(\bm{\overline x})| \leqslant \varepsilon 
\]
ce qui montre la continuité de $\phi$ est $\bm{\overline x}$ et donc par arbitrarité de $\bm{\overline x}$ on conclut que $\phi$ est continue sur $U$. \\\\
\textit{Continuité de $\frac{\partial \phi}{\partial x_i}$ : } Soient $\bm x_1 \neq \bm x_2$ dans $U$, $y_1 = \phi(\bm x_1)$ et	$y_2 = \phi(\bm x_2)$. Puisque $f \in C^1$, on a 
\[
0 = f(\bm x_2, y_1,) - f(\bm x_1, y_1) = \nabla f(\bm \xi, \eta) \cdot (\bm x_2 - \bm x_1, y_2 - y_1)
\]
Si l'on note $\nabla_{\bm x} f$ les $n$ première composantes on a,
\[
\nabla_{\bm x} f(\bm \xi, \eta) \cdot (\bm x_2 - \bm x_1) + \frac{\partial f}{\partial y} (\bm \xi, \eta)(y_2 - y_1) = 0
\]
Alors,
\[
\frac{\phi(\bm x_2) - \phi(\bm x_1)}{\norme{\bm x_2 - \bm x_1}} = - \frac{\nabla_{\bm x} f(\bm \xi, \eta)}{\frac{\partial f}{\partial y} (\bm \xi, \eta)} \cdot \frac{\bm x_2 - \bm x_1}{\norme{\bm x_2 - \bm x_1}}
\]
puisque $\frac{\partial f}{\partial y} > 0$ car $(\bm \xi, \eta) \in V \subset W$. On obtient donc par passage à la limite,
\[
\phi(\bm x_2) = \phi(\bm x_1) - \frac{\nabla_{\bm x} f(\bm x_1, y_1)}{\frac{\partial f}{\partial y}(\bm x_1, y_1)} \cdot (\bm x_2 - \bm x_1) + o(\norme{\bm x_2 - \bm x_1})
\]
donc $\phi$ est différentiable et par unicité de la différentielle,
\[
\forall i \in \entiers, \quad \frac{\partial \phi}{\partial x_i}(\bm x) = -\frac{\frac{\partial f}{\partial x_i} f(\bm x_1, y_1)}{\frac{\partial f}{\partial y}(\bm x_1, y_1)} 
\]
De plus les dérivées partielles de $\phi$ sont continues car celles de $f$ le sont. La démonstration que $\phi \in C^k$ se fait par récurrence sur $k \geqslant 1$.
}

\newcor{Si on considère l'hyperplan tangent à $\Sigma$ en $\bm z_0$, noté $\Pi_{\bm z_0} (\Sigma)$, comme l'hyperplan est tangent au graphe de $\phi$ en $\bm x_0$, qui est donnée par
	\[
	\Pi_{\bm z_0} (\Sigma) =  \left\{(\bm x, y) \in \R^{n+1} ~|~ y = \phi(\bm x_0) + \nabla \phi(\bm x_0) \cdot (\bm x - \bm x_0)\right\}
	\]
	Mais par le théorème précédent, l'hyperplan tangent peut s'écrire de façon équivalente comme,
	\[
	\frac{\partial f}{\partial y} (\bm z_0) (y - y_0) + \nabla_{\bm x} f (\bm z_0) \cdot (\bm x - \bm x_0) = 0
	\]
	Ce qui conclut à l'expression simple,
	\[
	\Pi_{\bm z_0}(\Sigma) = \{\bm z \in \R^{n+1} ~|~ \nabla f(\bm z_0) \cdot (\bm z - \bm z_0) = 0\}
	\]
}

\newthm{Soit $E \subset \R^{n+m}$ un ouvert non-vide, $\bm f \in C^1 (E, \rmm)$, $\Sigma = \{ (\bm x, \bm y) \in E ~|~ \bm f(\bm x, \bm y) = \bm 0\}$ et $\bm z_0 = (\bm x_0, \bm y_0) \in \Sigma$ tel que $\det D_{\bm y} \bm f(\bm z_0) \neq 0$. Alors il existe une boule ouvert $U = B(\bm x_0, \delta)$, un ouvert $V \subset E$ contenant $\bm z_0$ et une unique fonction $\bm \phi : U \to \rmm$ de classe $C^1$ tels que : 
	\begin{enumerate}
		\item $\bm y_0 = \bm \phi (\bm x_0)$
		\item $\forall \bm x \in U, \quad (\bm x, \bm \phi (\bm x)) \in U \et \bm f(\bm x, \bm \phi(\bm x)) = 0$
		\item $G(\bm \phi) = \Sigma \cap V$
	\end{enumerate}
	De plus, on a $\forall \bm x \in U, ~\det D_{\bm y} \bm f(\bm x, \bm \phi(\bm x)) \neq 0$ et,
	\[
	D \bm \phi(\bm x) = - D_{\bm y} \bm f(\bm x, \bm \phi(\bm x))^{-1} D_{\bm x} \bm f(\bm x, \bm \phi(\bm x))
	\]
	De plus si $\bm f \in C^k$ alors $\bm \phi C^k$.
}{Théorème des fonctions implicites - cas vectorielle}

\prf{Soit $\bm F : E \to \R^{n +m}$ définie par $\bm F(\bm x, \bm y) = (\bm x, \bm f(\bm x, \bm y))$. En particulier $\bm F(\bm x_0, \bm y_0) = (\bm x_0, \bm 0)$ et 
	\[
	D \bm F (\bm x_0, \bm y_0) = \begin{bmatrix}
		I_{n\times n} & 0 \\ D_{\bm x} \bm f(\bm z_0) & D_{\bm y} \bm f(\bm z_0)
	\end{bmatrix}
	,
	\quad 
	\det D \bm F (\bm x_0, \bm y_0) = \det D_{\bm y} \bm f(\bm z_0) \neq 0
	\]
	Donc il existe un ouvert $V' \subset E$ contenant $(\bm x_0, \bm y_0)$ et un ouvert $U'$ contenant $(\bm x_0, \bm 0)$ tels que $\bm F$ est une difféomorphisme de $V'$ à $U'$. Soit $\bm G$ la fonction inverse. On peut toujours trouver une boule ouverte $U = B(\bm x_0, \delta) \subset \rn$ et un ouvert $W \subset \rmm$ contenant $\bm 0$ tels que $U \times W \subset U'$. On considère alors la restriction $\bm F : V \to U \times W$ où $V = \bm G(U \times W)$. \\
	La fonction inverse à la forme $\bm G(\bm x, \bm w) = (\bm x, \bm \varphi(\bm x, \bm w))$ avec $\bm \varphi(\bm x_0, \bm 0) = \bm y_0$. On note, en particulier, que si $(\bm x, \bm y) \in \Sigma \cap V$, alors $\bm F(\bm x, \bm y) = (\bm x, \bm 0)$ et $\bm x \in U$. En prenant la fonction inverse on a $\bm G(\bm x, \bm 0) = (\bm x, \bm \varphi(\bm x, \bm 0)) = (\bm x, \bm y)$. La fonction implicite cherchée est alors $\phi (\bm x) = \bm \varphi(\bm x, \bm 0) : U \to \rmm$. En effet :
	\[
	\forall \bm x \in U, \quad (\bm x, \bm 0) = \bm F \circ \bm G(\bm x, \bm 0) = \bm F(\bm x, \bm \varphi(\bm x, \bm 0)) = \bm F(\bm x, \bm \phi(\bm x)) = (\bm x, \bm f(\bm x, \bm \phi(\bm x))) 
	\]
	donc pour tout $\bm x \in U, \quad \bm f(\bm x, \bm \phi(\bm x)) = 0$ et,
	\[
	\forall (\bm x, \bm y) \in \Sigma \cap V, \quad (\bm x, \bm y) \in \bm G \circ F (\bm x, \bm y) = \bm G(\bm x, \bm f(\bm x, \bm y)) = (\bm x, \bm \varphi(\bm x, \bm 0)) = (\bm x, \bm \phi(\bm x))
	\]
	donc $\Sigma \cap V = G(\bm \phi)$. \\
	De plus si $ \bm{\hat \phi} : U \to \rmm$ satisfait aussi $\Sigma \cap V = G(\bm{\hat \phi})$, alors $G(\bm \phi) = G(\bm{\hat \phi})$ et donc $\bm{\hat \phi}= \bm \phi$.\\
	Finalement, $\bm \phi$ est de classe $C^1$ puisque $\bm G$ est de classe $C^1$ car $\bm F$ est un difféomorphisme. De plus, $0 \neq \det D \bm F(\bm x, \bm y) = \det D_{\bm y} \bm f (\bm x, \bm y)$ pour tout $(\bm x, \bm y) \in v$, donc $\bm D_{\bm y} \bm f$ est inversible sur $V$. Par la formule dérivation de fonction composées on a,
	\[
	0 = D \bm f(\bm x, \bm \phi(\bm x)) = D_{\bm x} \bm f(\bm x, \bm \phi (\bm x)) + D_{\bm y} \bm f(\bm y, \bm \phi(\bm x)) \cdot D \bm \phi(\bm x)
	\]
	d'où
	\[
	D \bm \phi(\bm x) = - D_{\bm y} \bm f(\bm x, \bm \phi(\bm x))^{-1} D_{\bm x} \bm f(\bm x, \bm \phi(\bm x))
	\]
}

